\section{Evaluation}
\label{sec:evaluation}

This section reports the empirical evaluation of CloudAware-MOCCA using a
physical Android device and containerized edge and cloud services. The complete
dataset comprises 2{,}184 end-to-end task executions collected on physical
hardware over a Wi-Fi network: 1{,}928 from the rule-based and forced-baseline
sessions (which include a large targeted-mobility augmentation described in
Section~\ref{sec:decision-quality}), and 256 from a session in which the Random
Forest served as the live on-device planner.

\textbf{A note on dataset composition.}
The dataset was expanded beyond the original healthy-usage collection through a
targeted augmentation that forced the mobility branch (stationary, walking, and
vehicular states via a debug override) to enrich the under-represented cloud
class from 22 to 600 observations. This deliberately trades representativeness of
natural usage for per-class classifier coverage: the cloud share reported below
(32.5\%) is an artefact of that augmentation rather than the natural routing rate
of the policy, which is approximately 4\% under stationary conditions. Wherever a
result depends on the dataset reflecting natural usage, most importantly the
energy validation and the learned-versus-rules comparison, this dependence is
stated explicitly.


\subsection{Experimental Setup and Data Collection}
\label{sec:eval-setup}

\textbf{Testbed.}
Experiments were conducted using a Samsung Galaxy A56 running Android 14 and
connected over Wi-Fi to two Docker services hosted on a separate machine. The
edge service was constrained to 2~vCPU and 2~GB of memory, whereas the cloud
service was allocated 4~vCPU and 4~GB. Because both services were physically one
network hop from the device, a persistent 80~ms one-way network delay was
introduced on the cloud service to emulate wide-area communication latency. This
delay was an experimental control parameter rather than a measurement obtained
from a public cloud provider.

Five workloads spanning different levels of computational intensity were
evaluated: echo, SHA-256, grayscale conversion, matrix multiplication, and
video-frame edge detection. Echo and SHA-256 represent lightweight workloads,
grayscale conversion an intermediate workload, and matrix multiplication and
video-frame edge detection computationally intensive workloads. Payload size was
additionally varied for three workloads.

\textbf{Sessions.}
Eleven data-collection sessions were conducted under distinct and independently
verified conditions: a healthy baseline, a battery-discharge sweep, a four-level
network-degradation sweep applied simultaneously to both remote tiers, an offline
period, a processor-stress period, forced local-only and cloud-only baselines, a
mobility sweep, a payload-size sweep, and an edge-contention experiment involving
eight synthetic clients. Network degradation ranged from 100~ms with no packet
loss to 1000~ms with 30\% packet loss, while the additional cloud delay remained
in place. A subsequent targeted-mobility augmentation repeated the healthy,
mobility, battery, processor-stress, and payload blocks several times to grow the
dataset to 1{,}844 adaptive rows, principally to populate the cloud class and the
condition groups used in the regret analysis (Section~\ref{sec:decision-quality}).
Every dataset row corresponds to an actual task execution against the deployed
services; the augmentation rows differ only in that their movement state was
forced, and they carry the same override flag so that they can be excluded from
cost-model validation.

In addition, a further session executed the Random Forest policy as the live
on-device planner for 256 tasks (Section~\ref{sec:decision-quality}). Analyses
specific to the rule engine (Section~\ref{sec:cost-validation} and the
reliability figures below) are computed on the 1{,}928 rule-based and baseline
rows so that the learned-policy session does not perturb them, whereas the policy
comparisons of latency (Section~\ref{sec:system-performance}) and regret
(Section~\ref{sec:decision-quality}) treat the learned policy as an additional
policy alongside the rule engine and the two fixed baselines.

Because network degradation was introduced on the remote host, a manipulation
check was required to confirm that the imposed conditions were reflected in the
measurements observed by the mobile policy. Figure~\ref{fig:network-score-rtt}
plots the measured RTT to the edge tier against the corresponding network score,
excluding offline observations. The relationship is monotonically decreasing on a
logarithmic scale, and observations classified as unstable network conditions are
concentrated below the 0.30 decision threshold, confirming that the imposed
degradation was visible to the mobile policy rather than remaining a server-side
artefact.

\begin{figure}[t]
\centering
\includegraphics[width=\columnwidth]{network-score-vs-rtt.png}
\caption{Measured RTT versus computed network score.}
\label{fig:network-score-rtt}
\end{figure}


\subsection{System Performance: Reliability and Latency}
\label{sec:system-performance}

The overall fallback rate was \textbf{1.3\% (25/1{,}928)}, but the reliability
cost is not spread thinly. Every one of the 25 fallbacks (17 video-frame
edge-detection and 8 matrix-multiplication executions, all large-payload heavy
tasks) waited out the full client timeout of approximately 10.2~s (median
10{,}177~ms) before re-executing locally in a median of 35~ms. A fallback
therefore imposes a multi-second stall on the order of 290 times the local
execution it ultimately performs, so the aggregate rate understates the tail
experience, and the failures concentrate on precisely the workload the policy
should arguably not have offloaded (Section~\ref{sec:cost-validation}).
Independent server-side overload handling forwarded 18 of 946 remote requests
(1.9\%) from the edge tier to the cloud tier, recorded separately from direct
edge execution. The cloud-only baseline, for which local fallback was disabled,
completed all 40 requests successfully; this indicates that the corresponding
session did not encounter remote failure and should not be generalized to
environments with lower network reliability.

\begin{table}[t]
\centering
\caption{End-to-end latency by execution policy}
\label{tab:eval-policy-latency}
\footnotesize
\begin{tabular}{lrrrr}
\toprule
Policy & $n$ & Mean (ms) & Median (ms) & p95 (ms) \\
\midrule
Adaptive rule policy & 1819 & 708.7 & \textbf{61.0} & 6087.2 \\
Adaptive ML policy & 256 & 646.0 & \textbf{50.0} & 5547.0 \\
Local-only baseline & 44 & 66.6 & 39.0 & 181.0 \\
Cloud-only baseline & 40 & 1367.1 & 394.0 & 6852.9 \\
\bottomrule
\end{tabular}
\end{table}

Figure~\ref{fig:latency-policy} reports latency by execution policy, with mean
values and 95\% confidence intervals shown alongside p95 latency. Because the
distributions are strongly right-skewed, median latency is the load-bearing
statistic. The adaptive policy's value is best read against the two fixed
baselines: it delivers an approximately sixfold median-latency improvement over
always-cloud (61~ms versus 394~ms) while remaining within tens of milliseconds of
the always-local baseline (39~ms) on the light tasks that dominate the median.
The learned Random Forest policy tracks the rule engine closely (median 50~ms
versus 61~ms).

The direction of these gaps was confirmed with a percentile bootstrap (10{,}000
resamples) of the per-task difference in mean latency. The rule-based policy is
significantly faster than always-cloud on all five tasks, with every 95\%
confidence interval excluding zero (for example, video-frame edge detection
$-1668$~ms, CI $[-2435, -954]$). It is, however, significantly \emph{slower} than
always-local on all five tasks, most starkly for video-frame edge detection
($+5006$~ms, CI $[+4644, +5364]$), because local execution of that task takes a
median of only 170~ms whereas offloading its approximately 2.8~MB payload takes
roughly 5.8~s. For this workload and network, local execution is therefore often
the latency-optimal choice, and the adaptive policy's contribution is primarily
to avoid the always-cloud penalty rather than to beat local execution; a clear
latency win from offloading did not materialize for these tasks at these payload
sizes.

\begin{figure}[t]
\centering
\includegraphics[width=\columnwidth]{latency-by-policy.png}
\caption{End-to-end latency by execution policy.}
\label{fig:latency-policy}
\end{figure}


\subsection{Cost-Model Validation: Latency and Energy Estimation}
\label{sec:cost-validation}

Estimator predictions were compared with measured wall-clock latency using
1{,}354 trusted observations after excluding controlled-context and offline
sentinel records. Correlations across the three execution tiers ranged from 0.62
to 0.73, and remote estimates systematically underestimated observed latency by
approximately 360 to 870~ms. The estimator therefore captures relative
performance tendencies but is not yet calibrated for accurate absolute
prediction.

Whole-device energy measurements were also compared with the analytical model,
and the dataset expansion materially affected what could be validated. First, the
local energy coefficient can no longer be validated on this dataset: reliable
integration of battery current requires an execution of at least 500~ms, but
after the augmentation every local execution completed in under 223~ms (median
35~ms across 680 observations), because the heavier workloads that previously ran
locally were routed remotely under the forced-mobility conditions. No local
observation now satisfies the threshold, so the previously reported local
over-estimate of approximately $68\times$ is unvalidated on the present dataset,
a coverage regression introduced by the augmentation. Second, the remote
coefficient is accurate only for particular workloads: across 163 remote
observations of at least 500~ms the median ratio of measured to modeled energy
was 0.29, but this aggregate is dominated by the video-frame edge-detection
workload, for which the model over-estimates by roughly a factor of four, while
the few lightweight remote observations under-estimate by a factor of two to
four. The previously reported agreement within approximately 8\% therefore held
only under the original, lighter workload mix and is composition-dependent.

Beyond calibration, the remote energy model does not scale whole-device
background consumption with the complete round-trip duration; it prices only the
radio's transmit and idle draw. Figure~\ref{fig:energy-overrun} shows the ratio
of measured whole-device energy to the modeled local energy estimate against
remote round-trip time, excluding observations below 500~ms as sampling noise.
Among the 230 qualifying remote observations, 89 (39\%) consumed more whole-device
energy than the modeled local alternative, and the effect is concentrated in
video-frame edge-detection round trips (3.6 to 9.9~s at an approximately 2.8~MB
payload), the largest exceeding break-even by roughly a factor of 35. Combined
with the latency result in Section~\ref{sec:system-performance}, where offloading
that same task cost approximately 5.8~s against a 170~ms local execution, this
identifies a single root cause: both the latency and energy estimators under-price
the transmission of large payloads, so the policy offloads a task that is faster
and cheaper to run locally. That one mis-routing drives the latency tail (p95
6087~ms), the energy overruns, and all 25 fallbacks
(Section~\ref{sec:system-performance}). Scaling both the latency and idle-energy
terms by payload size and round-trip time, rather than by the radio's active
window alone, would remove it.

\begin{figure}[t]
\centering
\includegraphics[width=\columnwidth]{energy-overrun.png}
\caption{Remote energy overrun versus round-trip time.}
\label{fig:energy-overrun}
\end{figure}

Because energy is half of the cost model's objective, the whole-device energy
actually consumed under each policy is the directly relevant measure.
Table~\ref{tab:energy-policy} reports the median measured energy per task. The
learned-policy session is excluded because its device-power telemetry reads a
median of 2874~mW against approximately 30~mW in every other session, an
order-of-magnitude, session-specific artefact that makes its energy values
non-comparable. Local execution is the cheapest option for every task, because
the device returns to idle quickly rather than remaining awake for a network
round trip. The adaptive rule policy tracks local energy on the light tasks it
keeps on-device (1 to 6~mJ) but tracks cloud energy on the heavy tasks it
offloads (187~mJ for video-frame edge detection). The adaptive policy therefore
does not save whole-device energy relative to always-local; it is more efficient
than always-cloud, which pays the round-trip cost even for trivial tasks, but
energy is minimized by local execution, consistent with the overrun in
Figure~\ref{fig:energy-overrun}.

\begin{table}[t]
\centering
\caption{Median measured whole-device energy (mJ) per task, by policy}
\label{tab:energy-policy}
\footnotesize
\begin{tabular}{lrrr}
\toprule
Task & Local-only & Cloud-only & Rule-based \\
\midrule
echo & 1 & 9 & 2 \\
SHA-256 & 1 & 8 & 1 \\
grayscale & 1 & 13 & 5 \\
matrix-multiply & 1 & 10 & 6 \\
video-frame-edges & 5 & 192 & 187 \\
\bottomrule
\end{tabular}
\end{table}


\subsection{Decision Quality: Classifier and Regret Analysis}
\label{sec:decision-quality}

The 50-tree Random Forest trained on the twelve deployed features achieved a test
accuracy of 0.989 and a cross-validation accuracy of $0.999 \pm 0.002$, against a
majority-class baseline of 0.495. This quantifies fidelity to the deterministic
policy rather than independent decision quality, because the training labels were
generated by that policy. The test accuracy is lower than the 1.000 obtained on
the smaller pre-augmentation dataset; with 461 test observations a small number of
genuine boundary cases now appear, which is more defensible than perfect
reproduction. An alternative configuration incorporating three energy-related
variables slightly reduced both test and cross-validation accuracy (0.989 to
0.987 and 0.999 to 0.994), despite these features accounting for 14.5\% of Gini
importance, indicating that they were correlated with information already
represented by the deployed feature set.

The adaptive training subset contained 913 local, 331 edge, and 600 cloud
decisions. The 32.5\% cloud share is not the natural routing rate of the policy
but a consequence of the targeted-mobility augmentation; under naturally observed
stationary conditions the cloud share is approximately 4\%. The augmentation's
value is that cloud-tier precision and recall (both approximately 0.99 on 150
test observations) can now be reported at all, whereas the pre-augmentation
dataset provided only about six cloud test observations.

The twelve deployed features, ordered by their contribution to the trained
forest, are summarized in Table~\ref{tab:feature-importance}. The model relies
overwhelmingly on is\_stable and speedup, followed by the two latency estimates.
The prominence of is\_stable must be read in light of the augmentation, which
forced the mobility states so extensively that this feature became a
near-deterministic edge-versus-cloud discriminator; on the original dataset it
ranked below speedup and the latency estimates.

\begin{table}[t]
\centering
\caption{Deployed Random Forest features, ordered by contribution to the
trained model}
\label{tab:feature-importance}
\footnotesize
\begin{tabularx}{\columnwidth}{@{}lX@{}}
\toprule
Feature & Description \\
\midrule
is\_stable & Boolean; true when the device is stationary (smoothed accelerometer linear acceleration below 0.3~m/s$^2$). Gates edge eligibility: edge is selected only while stationary, cloud while moving. \\
speedup & Ratio of estimated local to estimated remote latency; the value on which the $1.5\times$ offload floor is applied. \\
est\_remote\_ms & Estimated remote execution latency, including transmission (ms). \\
est\_local\_ms & Estimated on-device execution latency (ms). \\
task\_complexity & Ordinal workload weight: light $=0$, medium $=1$, heavy $=2$. \\
network\_score & Composite network-health score in $[0,1]$ from RTT, bandwidth, and loss. \\
rtt\_ms & Measured round-trip time to the edge tier (ms). \\
battery\_percent & Device battery level (0 to 100\%). \\
cpu\_percent & Current device processor utilization (\%). \\
bandwidth\_mbps & Static per-transport downstream bandwidth estimate, not a live measurement. \\
is\_charging & Boolean; true when the device is charging. \\
network\_type\_rank & Ordinal connection type: none $=0$, LTE $=1$, Wi-Fi $=2$, 5G $=3$. \\
\bottomrule
\end{tabularx}
\end{table}

Classification accuracy characterizes the model offline but not as the active
planner. A dedicated session therefore executed the Random Forest as the live
on-device decision mechanism for 256 tasks with no fallback and no error,
producing a tier distribution of 55.1\% local, 40.2\% edge, and 4.7\% cloud and a
median latency of 50~ms. A direct comparison against the full rule-based pool is
confounded by the augmentation, which skewed that pool toward cloud (aggregate
50.2\% local, 16.3\% edge, 33.5\% cloud). Restricted to the matched stationary
conditions the learned-policy session actually covered, the rule engine routes
51.2\% local, 44.7\% edge, and 4.1\% cloud, so the cloud share realigns (4.1\%
versus 4.7\%) while the local and edge split differs modestly. The defensible
runtime-fidelity claim is therefore the narrower one: under matched conditions the
learned policy reproduces the cloud-gating behavior and the median latency of the
rule engine (50~ms versus 61~ms) at an inference cost below 1~ms per decision.

Because agreement with rule-generated labels cannot independently validate the
policy that generated them, decision quality was additionally assessed using
regret relative to a matched-condition empirical oracle. For each combination of
task, network, battery, and processor state with observations from at least two
tiers and at least three samples per tier, the tier with the lowest empirically
observed latency was treated as the reference. Twelve of the thirty-six populated
condition groups qualified once the augmentation observations were included.

\begin{table}[t]
\centering
\caption{Decision regret relative to the matched-condition empirical oracle}
\label{tab:regret}
\footnotesize
\begin{tabular}{lrrr}
\toprule
Policy & $n$ & Median regret (ms) & Best tier (\%) \\
\midrule
Local-only & 44 & 4.14 & 100.0 \\
Adaptive rule policy & 1422 & 40.62 & 45.0 \\
Adaptive ML policy & 103 & $-0.17$ & 28.0 \\
Static heuristic & 1609 & 123.92 & 40.6 \\
Cloud-only & 40 & 361.37 & 0.0 \\
\bottomrule
\end{tabular}
\end{table}

The static heuristic is a non-adaptive reference stronger than the fixed
baselines: it routes light tasks to local and medium or heavy tasks to edge when
stationary or cloud when moving, the policy's structure without any runtime cost
estimation. Scored on the same buckets, it reaches 123.92~ms median regret
against the rule engine's 40.62~ms, so the adaptive cost model yields roughly a
threefold reduction in median regret over a plausible static rule of thumb. This
is not a re-implementation of a specific published algorithm (a direct comparison
against prior offloading policies remains future work,
Section~\ref{sec:sensitivity-summary}), but it establishes that adaptivity earns
its keep against more than the trivial baselines.

Figure~\ref{fig:regret-policy} shows regret with 95\% confidence intervals. The
favorable result for local-only reflects the composition of the qualifying
condition groups and should not be read as general superiority. Both adaptive
policies avoid the worst-tier choice made by the cloud-only baseline. Two
properties sit in tension and are both reported: the learned policy attains the
lowest median regret of any policy ($-0.17$~ms, indicating selections on average
as fast as the mean of the oracle-best tier), yet it selects the exact best tier
less often than the rule engine (28.0\% versus 45.0\%). This is consistent with a
smoothed classifier that avoids large errors but is less precise at decision
boundaries. Because the estimate rests on twelve small groups, the best-tier
percentages are noisy and should not be over-interpreted; the learned-policy
session also covered fewer conditions than the rule-based collection and predates
the augmentation. The defensible conclusion is that both adaptive policies
outperform the fixed cloud-only policy on this label-independent metric, and that
the runtime quality of the learned policy is at least comparable to that of the
rule engine.

\begin{figure}[t]
\centering
\includegraphics[width=\columnwidth]{regret-by-policy.png}
\caption{Decision regret relative to the matched-condition oracle.}
\label{fig:regret-policy}
\end{figure}


\subsection{Sensitivity and Threats to Validity}
\label{sec:sensitivity-summary}

\textbf{Mobility.}
Because the physical device remained stationary during laboratory collection,
mobility states were controlled experimentally, extensively so during the
augmentation. Figure~\ref{fig:mobility-routing} shows target-tier selection under
these conditions for adaptive executions. The edge tier was not selected in any
of the 1{,}162 observations recorded while the device was moving, confirming
enforcement of the stationary-only edge eligibility condition across a large
number of opportunities to violate it.

\begin{figure}[t]
\centering
\includegraphics[width=\columnwidth]{mobility-routing.png}
\caption{Target-tier selection under controlled mobility conditions.}
\label{fig:mobility-routing}
\end{figure}

\textbf{Payload size.}
For matrix multiplication, median latencies were 222~ms at the edge tier and
98~ms at the cloud tier with negligible correlation against payload size, and the
SHA-256 and grayscale workloads were similarly flat (absolute correlation at most
0.23). No clear relationship between payload size and transmission cost was
observed over the tested range.

\textbf{Threats to validity.}
The evaluation used a single device and one network, so device-specific
quantities should not be treated as universal constants. The 80~ms cloud delay
was a chosen emulation parameter, not a datacenter measurement. Edge contention
was generated with synthetic HTTP clients on the experimental host rather than
independent devices, so it demonstrates responsiveness to server load rather than
multi-user scalability. Most importantly, the cloud class was enriched by forcing
the mobility branch, so the dataset over-represents moving-device conditions; the
classifier's per-class metrics benefit, but the class distribution, the loss of
local-energy coverage, and the learned-versus-rules comparison are all affected.
The regret analysis covers twelve of thirty-six groups on small samples, and
bandwidth is a static capability estimate rather than measured throughput.

\textbf{Not yet measured (future work).}
Four questions this evaluation does not answer are stated so they are not mistaken
for settled. First, middleware overhead: the continuous cost of the MAPE loop,
the context collectors, and inference in CPU and battery terms is unquantified
beyond the sub-millisecond per-decision figure, and an idle-drain measurement is
required to confirm the system is not a net energy loss. Second, adaptation
latency: the time between a context change and the policy acting on it, the
defining property of a self-adaptive system, is not measured. Third, WAN
sensitivity: the edge-versus-cloud trade-off pivots on the chosen 80~ms cloud
delay, which was not swept, so an edge-favoring outcome could be an artefact of
that constant. Fourth, comparison to prior work: the static heuristic is the only
non-trivial baseline, and re-implementing an established offloading policy would
position the results in the literature. The first three require additional
on-device collection; the fourth requires implementing a baseline algorithm.

\textbf{Summary.}
The adaptive policy's demonstrated value is defensive: it delivers an
approximately sixfold median-latency improvement over always-cloud (61~ms versus
394~ms) at a 1.3\% fallback rate, avoids the worst-tier choice a fixed cloud-only
policy makes on the non-circular regret metric, and beats a static complexity
heuristic by roughly threefold in median regret. What it does not do, for this
workload and network, is beat local execution: always-local was significantly
faster on every task and consumed the least whole-device energy
(Table~\ref{tab:energy-policy}), so the policy's role here is largely to recognize
when to stay local. The Random Forest reproduces the rule engine with high
fidelity (0.989 test, 0.999 cross-validation accuracy), imitation fidelity rather
than an independent quality signal, and a 256-decision runtime session confirms
the imitation holds on matched conditions. The clearest actionable finding is a
cost-model defect: the latency and energy estimators under-price large-payload
transmission, so the policy offloads video-frame edge detection (170~ms local
versus approximately 5.8~s offloaded), and that single mis-routing accounts for
the latency tail, the 39\% remote-energy overrun, and all 25 fallbacks.
Correcting the transmission term is the highest-value next step.
